% Proposed new subsubsection for Chapter_03_Validation/Solid/Ablative/AblatingSolid.tex,
% to go after "Two-Phase Stefan Problem" (sec:val_abl_stefan2), replacing the WIP sentence on the
% surface that starts below the ablation temperature.  Figures to copy into
% Chapter_03_Validation/figures/Ablation/ from flames/tests/FlowAblation/reference/Images/R32_Onset/.
% Deck: tests/FlowAblation/input_1D_AblationOnset; reference: reference/ablation_onset_reference.py;
% check: reference/ablation_onset_check.py.  Runs: bin/tests/FlowAblation/R32_Onset_N100sq, _N200, _N400p.

\textbf{CLAUDE-START}
%-------------------------------------------------------------
\subsubsection{Onset of Ablation}
\label{sec:val_abl_onset}

In the two previous tests the surface is either never at the ablation temperature (Section~\ref{sec:val_abl_contact}) or at it from the start (Section~\ref{sec:val_abl_stefan2}). Here the surface starts \emph{below} $\Tabl$ and reaches it later, so both regimes of the capped wall temperature, Eq.~\ref{eq:Tw_capped}, and the switch between them are exercised in one run. A solid slab $0 < x < L = 0.1$ with an adiabatic back face at $x = 0$ is brought into contact at $t = 0$ with gas at rest at $T_g = 1$ filling $0.1 < x < 0.5$ (adiabatic end). The gas is that of Section~\ref{sec:val_abl_stefan2} ($k_T = 4.861\times10^{-3}$, $\alpha_g = 1.389\times10^{-3}$). The solid has $\rho_s = 1$, $c_s = 0.875$ and $k_s = 4k_T$, chosen so that its thermal effusivity equals that of the gas, $e_s = e_g$, while its diffusivity is $\alpha_s = 16\,\alpha_g$, so the slab heats through on a time scale comparable to the gas diffusion time. The solid starts at $T_{s0} = 0.99$, $\Tabl = 0.9975$ and $\Qstar = 0.01$. The contact temperature of Eq.~\ref{eq:contact_temperature} is $T_i = 0.995 < \Tabl$: at first nothing ablates. As heat accumulates against the adiabatic back face, the surface temperature rises, reaches $\Tabl$ at $t_{abl}$, and from then on the surface is held at $\Tabl$ and the excess flux removes solid.

\paragraph{Reference solution.}
Before onset the problem is the contact of a semi-infinite gas with a finite slab. By Laplace transform (or the method of images in the slab~\cite{CarslawJaeger1959}) the surface temperature is
\begin{equation}
    T_w(t) = T_{s0} + \frac{T_g - T_{s0}}{1+\sigma}\left[1 + (1+\beta)\sum_{n\ge1}\beta^{\,n-1}\,\mathrm{erfc}\!\left(\frac{nL}{\sqrt{\alpha_s t}}\right)\right],
    \qquad \sigma = \frac{e_s}{e_g},\quad \beta = \frac{\sigma - 1}{\sigma + 1},
    \label{eq:onset_Tw}
\end{equation}
which starts at $T_i$ and tends to $T_g$. With $\sigma = 1$ the series reduces to a single term, and $T_w = \Tabl$ gives the onset time in closed form,
\begin{equation}
    t_{abl} = \frac{L^2}{\alpha_s\left[\mathrm{erfc}^{-1}\!\big(2(\Tabl - T_i)/(T_g - T_{s0})\big)\right]^2} = 1.97830 .
    \label{eq:onset_time}
\end{equation}
After onset the surface obeys $\rho_s\Qstar\,ds/dt = q_g - q_s$ with $T_w = \Tabl$, starting from the non-similar temperature field of the heating stage, and has no closed form. We compute it with a separate one-dimensional solver: each phase is mapped onto a fixed grid by a front-fixing (Landau) transformation, the heat equations are discretised with second-order central differences and second-order one-sided surface gradients, and the system is integrated implicitly with the switch between the two surface conditions located as an event. The reference solver uses the same finite gas column as the run. We verified it in three ways. First, its onset time converges at second order to Eq.~\ref{eq:onset_time} ($-0.071$, $-0.018$, $-0.0045\%$ on 200, 400, 800 solid nodes). Second, with the parameters of Section~\ref{sec:val_abl_stefan2} it reproduces the two-phase Stefan similarity solution ($+0.036$, $+0.009$, $+0.002\%$ in $s$ at $t = 1$ on 1000, 2000, 4000 nodes). Third, on the present problem it gives $s(t = 6) = 0.024085$, which changes by less than $10^{-5}$ in relative terms between the two finest grids. Like the previous tests, the reference neglects the motion and density change of the gas and assumes no blowing.

\paragraph{Results.}
The surface does not move before onset: the recession stays below $10^{-13}$ until the uncapped surface temperature reaches $\Tabl$. The numerical onset time is taken from the uncapped surface temperature of Eq.~\ref{eq:Tw_capped} at the last three plot times before the surface leaves its initial face, extrapolated to $\Tabl$. Table~\ref{tab:val_abl_onset} lists the errors on three grids. All are below $1\%$ and decrease with the cell size, at an observed order of about $0.6$ between the two finest grids. The onset time is ill-conditioned: near onset $dT_w/dt \approx 5.6\times10^{-4}$, so a surface-temperature error of $0.01\%$ of $\Delta T = T_g - T_{s0}$ shifts $t_{abl}$ by about $0.1\%$. Before onset the surface temperature agrees with Eq.~\ref{eq:onset_Tw} to within $0.15\%$ of $\Delta T$ for $t \geq 1$ and $0.07\%$ for $t \geq 1.5$ on all grids. The largest deviation, $0.5\%$ of $\Delta T$ at $t \approx 0.7$ on 400 cells, follows the gas pressure at the wall (correlation coefficient $0.98$). Its period matches the acoustic round trip $2 \times 0.4/c \approx 0.68$: the sudden cooling of the gas at contact launches pressure waves that reflect between the wall and the open end of the domain, which the conduction-only reference cannot contain. The gas temperature profile at $t = 6$ differs from the reference by $0.07$--$0.09\%$ of $\Delta T$, nearly independent of the grid. This is the gas-density effect already seen in Sections~\ref{sec:val_abl_contact} and~\ref{sec:val_abl_stefan2}.

\begin{table}[H]
    \centering
    \begin{tabular}{lccc}
        \toprule
        Cells ($\Delta x$) & Onset time $t_{abl}$ & Recession $s(t = 6)$ & Gas profile, $t = 6$ \\
        \midrule
        $100$ ($5.0\times10^{-3}$)  & $1.9843$ ($+0.30\%$) & $0.023829$ ($-1.07\%$) & $0.087\%$ of $\Delta T$ \\
        $200$ ($2.5\times10^{-3}$)  & $1.9828$ ($+0.23\%$) & $0.023906$ ($-0.75\%$) & $0.081\%$ of $\Delta T$ \\
        $400$ ($1.25\times10^{-3}$) & $1.9812$ ($+0.15\%$) & $0.023972$ ($-0.47\%$) & $0.070\%$ of $\Delta T$ \\
        \midrule
        Reference & $1.97830$ (Eq.~\ref{eq:onset_time}) & $0.024085$ & -- \\
        \bottomrule
    \end{tabular}
    \caption{Onset of ablation, Numerical against Reference: onset time (closed form, Eq.~\ref{eq:onset_time}), recession at $t = 6$ (front-fixing reference solution) and maximum temperature error in the gas at $t = 6$. The solid temperature error is below $10^{-3}\%$ of $\Delta T$ on all grids.}
    \label{tab:val_abl_onset}
\end{table}

\begin{figure}[H]
    \centering
    \includegraphics[width=1.0\linewidth]{Chapter_03_Validation/figures/Ablation/ablation_onset_N400.png}
    \caption{Onset of ablation (400 cells): surface temperature history, surface recession, temperature profile at $t = 6$ and its error, Numerical against Reference. The reference is Eq.~\ref{eq:onset_Tw} before onset and the front-fixing solution after it. The dashed line marks $\Tabl$ and the dotted line the exact onset time, Eq.~\ref{eq:onset_time}. The maximum errors are $0.15\%$ in the onset time, $0.47\%$ in the recession and $0.07\%$ of $\Delta T$ in the temperature.}
    \label{fig:val_abl_onset}
\end{figure}

\begin{figure}[H]
    \centering
    \includegraphics[width=0.75\linewidth]{Chapter_03_Validation/figures/Ablation/ablation_onset_convergence.png}
    \caption{Onset of ablation: grid convergence of the onset time, the recession at $t = 6$ and the temperature profile error. All errors are below $1.1\%$ on 100 cells and below $0.5\%$ on 400 cells.}
    \label{fig:val_abl_onset_conv}
\end{figure}
\textbf{CLAUDE-END}
